\section{Introduction}
\label{sec:intro}

\paragraph{The 99.4\,\% problem.}
Static-prefix multi-agent LLM coding systems assign dispatched
subagents to cheap-to-expensive tiers (haiku / sonnet / opus) via
a (role~\(\times\)~signal) rule cascade. When the cascade is
well-tuned---as it is in \blackrim, where 16 static rules route
each subagent role to its cost-optimal tier (Researcher to haiku,
Builder to sonnet, Architect to opus)---the companion paper
\cite{moa1landscape} reports an 8.9\(\times\) cost reduction on
subagent dispatch alone. Yet this improvement makes a different
problem more visible. Production telemetry measured over a 7-day
window of \(\sim\)50 spawns per day shows that total session cost
splits as \textbf{99.4\,\%} on the main-thread orchestrator
(Gestalt running on claude-opus-4-7) and \textbf{0.6\,\%} on all
dispatched subagents combined (Table~\ref{tab:cost-split},
\S\ref{sec:eval}). The implication is stark: once subagent
dispatch is tuned, further dispatch optimisation yields at most
0.6\,\% additional savings. The orchestrator itself is the
binding cost lever, and it has no analogous tier selector. Every
main-thread turn pays opus rates regardless of whether the turn
requires opus-level capability.

\paragraph{Two complementary primitives.}
Two architectural primitives target the orchestrator cost directly.
\emph{Per-turn model routing} selects the cheapest model tier
adequate for each main-thread turn---replacing the implicit
always-opus default without sacrificing the quality needed for
orchestration-level reasoning. The main-thread turn distribution
is heterogeneous: a bd-claim confirmation, a read-only status
question, and a multi-step architecture decision all appear on
the same thread but demand very different model capability levels.
A well-calibrated router can route the first two to haiku or
sonnet while reserving opus for the third.
\emph{Plan-level semantic caching} short-circuits the briefing-composition
step for structurally equivalent bd dispatches, stacking on top
of the existing prefix-cache layer rather than competing with it.
Neither primitive has a published methodology grounded in
production multi-agent telemetry: prior routing work targets
either subagent dispatch (where role labels are available) or
single-turn chat (where the conversation is a single pair, not
a long-running orchestration session), and prior caching work
treats prefix reuse within a fixed TTL window rather than
plan-level reuse across an unbounded horizon.

\paragraph{Our approach.}
We design and evaluate both primitives on \blackrim's production
instance. The \emph{semantic-similarity router} (\S\ref{sec:routing})
represents each main-thread turn as an embedding, finds its
nearest neighbour in a hand-curated exemplar set of
\(\sim\)150 labeled turns, and routes to the exemplar's tier---
escalating conservatively to opus when no exemplar exceeds a
similarity threshold~\(\tau\). This approach requires no training,
tolerates a small exemplar set, and produces interpretable routing
decisions. The \emph{agentic plan cache} (\S\ref{sec:plancache})
derives a stable semantic signature from the orchestration plan
(agent role, task shape, normalised file scope, deliverable kind),
performs exact-then-approximate lookup, and returns the cached
output on hit, short-circuiting the \texttt{gt context budget}
/ \texttt{gt orchestrate plan} pipeline entirely.

\paragraph{Contributions.}
This paper makes four contributions:

\begin{enumerate}[label=(\roman*),leftmargin=2em]
  \item \textbf{Cost-decomposition measurement.} We report the
    99.4\,\% / 0.6\,\% main-thread / dispatch cost split for a
    production multi-agent coding system and explain why it
    inverts the optimisation target relative to the prior routing
    literature.

  \item \textbf{Semantic-similarity per-turn router.} We implement
    and evaluate a k-NN router over sentence embeddings, with
    conservative escalation, on a 49-turn hand-labeled test set
    drawn from real session transcripts. Leave-one-out
    cross-validation yields macro-F1 0.526 (+6.0 points over the
    length-heuristic baseline), with haiku-F1 0.786 and sonnet-F1
    0.667. Estimated cost savings are 30\,\% above the
    length-heuristic baseline.

  \item \textbf{Plan-cache calibration measurement.} We report
    the existing prefix-cache baseline---10.44\(\times\)
    read-to-creation ratio, 90.1\,\% hit rate, 78.6\,\% estimated
    savings on real-API spawns from Anthropic's
    \texttt{cache\_control} layer---as the calibration point atop
    which plan-cache incremental savings are measured. We define
    the plan-signature function and describe the cache-fill and
    eviction policy.

  \item \textbf{Composed-cost analysis.} We show that per-turn
    routing and plan caching are independent in the cost equation
    and that their savings compound multiplicatively, enabling a
    single closed-form cost-reduction bound for a session that uses
    both.
\end{enumerate}

We also document an integration blocker---the current Claude Code
interface exposes session-level model selection only, not per-turn
routing---and evaluate three workaround paths
(\S\ref{sec:discussion}).

\paragraph{Paper structure.}
\S\ref{sec:related} surveys the two literature strands this paper
draws on: LLM model routing (cascading, preference-data routing,
confidence-based escalation) and caching for LLM systems (prefix
caching, agentic plan caching). \S\ref{sec:problem} formalises the
per-turn routing and plan-cache decisions as a joint cost-minimisation
problem. \S\ref{sec:system} describes \blackrim's relevant
architecture. \S\ref{sec:routing} presents the semantic-similarity
router. \S\ref{sec:plancache} presents the plan-cache design.
\S\ref{sec:eval} reports empirical evaluation. \S\ref{sec:discussion}
discusses integration constraints and limitations.
\S\ref{sec:future} maps future directions.
